% ================= SEKCJA 1 =================
\section{Tytuł}

Oficjalnym tytułem projektowanego rozwiązania jest: \textbf{WorkFlow -- system wspierający zarządzanie pracownikami, czasem pracy i dokumentacją kadrową}. Tytuł ten odzwierciedla główny zakres aplikacji, który obejmuje ewidencję pracowników, przypisywanie ich do zakładów i projektów, rejestrację czasu pracy, obsługę dokumentów kadrowych, kontrolę certyfikatów oraz przygotowanie danych do rozliczeń.

Nazwa \texttt{WorkFlow} została przyjęta jako nazwa techniczna i projektowa systemu. Wskazuje ona na uporządkowany przepływ informacji pomiędzy poszczególnymi modułami aplikacji oraz rolami użytkowników. W projekcie istotne jest nie tylko samo przechowywanie danych, lecz także zachowanie zależności pomiędzy pracownikiem, zakładem, projektem, czasem pracy, dokumentacją oraz procesem rozliczeniowym.

\subsection{Definicja i zakres przedmiotowy}

System \texttt{WorkFlow} jest aplikacją webową typu full-stack, zaprojektowaną w celu wsparcia podstawowych procesów kadrowo-organizacyjnych w przedsiębiorstwie. Rozwiązanie składa się z warstwy frontendowej, warstwy backendowej, relacyjnej bazy danych, magazynu plików oraz infrastruktury uruchomieniowej opartej o kontenery.

Zakres przedmiotowy projektu obejmuje kilka powiązanych ze sobą obszarów:
\begin{itemize}
    \item \textbf{Zarządzanie pracownikami} -- prowadzenie centralnej ewidencji pracowników, ich danych podstawowych, statusu aktywności, roli systemowej oraz przypisania do zakładu.
    \item \textbf{Obsługa zakładów i projektów} -- organizowanie danych według zakładów oraz przypisywanie pracowników do projektów, które reprezentują miejsca lub obszary wykonywania pracy.
    \item \textbf{Rejestracja czasu pracy} -- przyjmowanie zdarzeń wejścia i wyjścia z czytników RCP oraz tworzenie na ich podstawie wpisów czasu pracy.
    \item \textbf{Dokumentacja kadrowa} -- przechowywanie informacji o umowach, dokumentach, certyfikatach, badaniach, szkoleniach oraz nieobecnościach pracowników.
    \item \textbf{Rozliczenia} -- przygotowanie uporządkowanych danych na potrzeby rozliczeń kadrowo-płacowych, w szczególności na podstawie zatwierdzonych wpisów czasu pracy.
\end{itemize}

W takim ujęciu system nie jest pełnym systemem ERP ani kompletnym narzędziem księgowym. Jego zadaniem jest wsparcie wybranego zakresu procesów związanych z obsługą pracowników i czasu pracy oraz przygotowanie danych, które mogą zostać wykorzystane przez osoby odpowiedzialne za kadry i rozliczenia.

\subsection{Charakterystyka rozwiązania}

Projekt ma charakter praktyczny i został przygotowany jako działająca aplikacja, możliwa do uruchomienia w środowisku lokalnym lub serwerowym. Przyjęta architektura rozdziela odpowiedzialności pomiędzy kilka warstw systemu:
\begin{itemize}
    \item \textbf{frontend} odpowiada za interfejs użytkownika i komunikację z użytkownikiem końcowym,
    \item \textbf{backend} realizuje logikę biznesową, autoryzację, walidację danych oraz komunikację z bazą danych,
    \item \textbf{baza danych PostgreSQL} przechowuje dane relacyjne systemu,
    \item \textbf{MinIO} służy jako magazyn plików i dokumentów,
    \item \textbf{Nginx} pełni funkcję reverse proxy dla aplikacji.
\end{itemize}

Istotnym elementem projektu jest również podział użytkowników według ról. System przewiduje role takie jak administrator, pracownik HR, księgowość, brygadzista, pracownik oraz rola biurowa. Dzięki temu dostęp do funkcji i danych może być ograniczany zgodnie z zakresem odpowiedzialności danej osoby.

\subsection{Zakres dokumentacji}

Niniejsza dokumentacja przedstawia zarówno założenia projektowe, jak i techniczne aspekty wykonania systemu. Obejmuje ona opis celu projektu, analizę wymagań, przypadki użycia, scenariusze działania, makiety interfejsu, model danych oraz zastosowany stos technologiczny.

Szczególną uwagę poświęcono modelowi danych, ponieważ to on stanowi podstawę działania większości modułów systemu. Relacje pomiędzy pracownikami, zakładami, projektami, czasem pracy, dokumentami, certyfikatami i rozliczeniami zostały przedstawione w postaci diagramu ERD oraz opisane w dalszej części dokumentu.

\subsection{Znaczenie dla dalszej części dokumentu}

Tak sformułowany tytuł i zakres projektu stanowią punkt wyjścia dla dalszej analizy. W kolejnych rozdziałach opisano nazwę kodową systemu, cele jego realizacji, wymagania funkcjonalne i niefunkcjonalne, scenariusze użycia oraz strukturę danych.

Dokumentacja ma na celu przedstawienie aktualnego stanu projektu w sposób uporządkowany i możliwy do zweryfikowania. Z tego względu opis skupia się na funkcjach rzeczywiście przewidzianych lub zaimplementowanych w systemie, bez rozszerzania zakresu o moduły, które nie są częścią obecnej wersji aplikacji.